\section{Analysis}
\label{sec:analysis}

We analyze the runtime characteristics of ISD serving to understand where time is spent, why certain bottlenecks are fundamental to strided decoding, and what headroom remains for future optimization.

\subsection{Step Latency Decomposition}
\label{sec:step_breakdown}

Table~\ref{tab:step_breakdown} breaks down the per-step latency of ISD at stride $N{=}3$ on a single H100 with CUDA graph enabled, at concurrency $C{=}1$.

\begin{table}[t]
\centering
\caption{\textbf{Step latency breakdown} for ISD ($N{=}3$, $C{=}1$, 1$\times$H100). GPU forward dominates; CPU overhead is serial and directly additive.}
\label{tab:step_breakdown}
\small
\begin{tabular}{lrr}
\toprule
\textbf{Component} & \textbf{Latency (ms)} & \textbf{\% of Step} \\
\midrule
GPU forward (CUDA graph replay) & 7.06 & 83.1\% \\
Verify + sampling (Triton kernel) & 0.76 & 8.9\% \\
KV trim + output assembly & 0.25 & 2.9\% \\
Batch preparation (classify + metadata) & 0.16 & 1.9\% \\
Scheduling + misc overhead & 0.27 & 3.2\% \\
\midrule
\textbf{Total step time} & \textbf{8.50} & 100\% \\
\bottomrule
\end{tabular}
\end{table}

The GPU forward pass (CUDA graph replay) accounts for the majority of step time. Importantly, all non-forward components run \emph{serially} on the CPU---unlike AR decoding, where CPU scheduling is hidden behind GPU compute via overlap scheduling. This serial dependency is intrinsic to ISD: the verify step requires the forward pass's logits, and the next step's batch preparation requires the verify step's accept/reject decisions. Every millisecond of CPU overhead is directly additive to wall-clock latency.

\subsection{The Serial Dependency Bottleneck}
\label{sec:serial_bottleneck}

The fundamental performance constraint of ISD is the strict dependency chain within each step:

\vspace{-4pt}
\begin{equation*}
\text{forward}(t) \;\to\; \text{verify}(t) \;\to\; \text{trim}(t) \;\to\; \text{prepare}(t{+}1) \;\to\; \text{forward}(t{+}1)
\end{equation*}
\vspace{-4pt}

AR serving systems break an analogous chain by exploiting the fact that the next step's batch preparation depends only on \emph{which} requests are active, not on the \emph{content} of the current step's output. This allows CPU scheduling to run on a separate stream, overlapping with GPU compute. ISD cannot use this technique: the verify step produces accept/reject decisions that determine how many KV slots to trim and where the next step's sequence begins. The prepare step must wait for these decisions.

This bottleneck manifests most clearly in multi-GPU scaling. Table~\ref{tab:tp_scaling} compares ISD and AR scaling from TP=1 to TP=4 on H100s.

\begin{table}[t]
\centering
\caption{\textbf{Multi-GPU scaling.} AR achieves near-linear speedup from TP=1 to TP=4 because CPU overhead is hidden. ISD's serial CPU overhead becomes a larger fraction of (reduced) GPU forward time, limiting scaling.}
\label{tab:tp_scaling}
\small
\begin{tabular}{lccc}
\toprule
& \textbf{TP=1 (tok/s)} & \textbf{TP=4 (tok/s)} & \textbf{Scaling} \\
\midrule
Qwen3-8B AR ($C{=}1$) & --- & --- & ---$\times$ \\
Ours $N{=}3$ ($C{=}1$) & 285 & --- & ---$\times$ \\
\bottomrule
\end{tabular}
\end{table}

With tensor parallelism, the GPU forward time decreases (more SMs available), but the serial CPU overhead---verify, trim, batch preparation---remains constant. As GPU time shrinks, the fixed CPU overhead becomes a larger fraction of total step time, yielding sublinear scaling. At TP=4, the CPU overhead of ${\sim}$1.5ms constitutes roughly 30\% of the total step time, compared to $<$5\% for AR (where it is hidden by overlap).

\subsection{Theoretical Throughput Limits}
\label{sec:theoretical}

We can bound the maximum achievable throughput of ISD by separating hardware-limited and software-limited components.

\paragraph{Hardware limit.}
At stride $N$ with acceptance rate $\alpha$, the expected tokens per forward pass is $\text{TPF}_N$ (Eq.~\ref{eq:tpf}). If the GPU forward takes $T_{\text{fwd}}$ seconds and all other overhead is zero, the theoretical maximum throughput is:
\begin{equation}
\text{TPS}_{\max} = \frac{\text{TPF}_N}{T_{\text{fwd}}}.
\label{eq:tps_max}
\end{equation}

\paragraph{Current efficiency.}
Table~\ref{tab:efficiency} reports the ratio of measured throughput to the hardware-limited maximum across configurations, revealing how much headroom remains.

\begin{table}[t]
\centering
\caption{\textbf{Throughput efficiency} relative to the hardware-limited maximum (Eq.~\ref{eq:tps_max}). The gap is entirely due to serial CPU overhead.}
\label{tab:efficiency}
\small
\begin{tabular}{lcccc}
\toprule
\textbf{Config} & \textbf{TPF} & $T_{\text{fwd}}$ \textbf{(ms)} & \textbf{Measured TPS} & \textbf{Efficiency} \\
\midrule
$N{=}3$, TP=1, $C{=}1$ & 2.47 & 7.06 & 285 & 81.4\% \\
$N{=}5$, TP=1, $C{=}1$ & 2.71 & 7.11 & 303 & 79.5\% \\
$N{=}3$, TP=4, $C{=}1$ & \multicolumn{4}{c}{\emph{see future work}} \\
\bottomrule
\end{tabular}
\end{table}

The efficiency gap is attributable entirely to serial CPU overhead: verify, KV trim, batch preparation, and metadata updates. At TP=1, the GPU forward dominates and efficiency is relatively high. At TP=4, the forward time drops but CPU overhead is constant, reducing efficiency.

\subsection{Stride Size Trade-offs}
\label{sec:stride_tradeoff}

Increasing the stride $N$ has three competing effects on throughput:

\begin{denseitemize}
\item \textbf{Higher TPF}: More proposals per forward pass means more potential tokens accepted per step, directly increasing the numerator of TPS.
\item \textbf{Longer forward time}: The extend attention kernel processes $2N{-}1$ tokens instead of $2N'{-}1$, increasing $T_{\text{fwd}}$. At small batch sizes this effect is modest (memory-bound regime), but at high concurrency the extra tokens consume compute.
\item \textbf{Lower per-proposal acceptance}: Later proposals condition on earlier (potentially incorrect) proposals, compounding errors. The marginal acceptance rate decreases with position, yielding diminishing returns in TPF.
\end{denseitemize}

Table~\ref{tab:stride_ablation} quantifies these trade-offs.

\begin{table}[t]
\centering
\caption{\textbf{Stride size ablation} ($C{=}1$, 1$\times$H100). $N{=}3$ achieves the best throughput; larger strides have diminishing TPF gains that are offset by longer forward times.}
\label{tab:stride_ablation}
\small
\begin{tabular}{lcccc}
\toprule
\textbf{Stride $N$} & \textbf{TPF} & \textbf{Accept \%} & \textbf{TPS} & $T_{\text{step}}$ \textbf{(ms)} \\
\midrule
$N{=}3$ & 2.47 & 82.6\% & 285 & 8.5 \\
$N{=}4$ & 2.78 & 65.2\% & 301 & 8.5 \\
$N{=}5$ & 2.71 & 38.5\% & 303 & 8.6 \\
\bottomrule
\end{tabular}
\end{table}

\subsection{Future Directions}
\label{sec:future}

The analysis above identifies the serial CPU dependency chain as the primary remaining bottleneck. We outline three directions to address it:

\paragraph{Speculative metadata pre-computation.}
Since 78\% of proposals are accepted, the next step's sequence lengths can be predicted with high accuracy during the current GPU forward. Pre-computing attention metadata (flashinfer plans, page table updates) speculatively during GPU compute would hide the metadata cost in the common case, falling back to recomputation only on misprediction.

\paragraph{Hybrid decode--extend attention.}
The prefix portion of ISD attention is structurally identical to AR decode attention (one-to-many), while only the new $2N{-}1$ tokens require extend-style (many-to-many with causal masking). Splitting the attention into a decode kernel for the prefix and a local causal kernel for new tokens would reduce metadata cost from ${\sim}$0.43ms to ${\sim}$0.18ms, since the decode metadata path is highly optimized in existing systems.

\paragraph{Pipelined verification.}
Breaking the strict dependency chain by pipelining verify$(t)$ with forward$(t{+}1)$---launching the next forward pass speculatively before verification completes---would recover the CPU--GPU overlap that AR systems enjoy. The pipeline would need rollback support for the rare case (${\sim}$22\%) where verification invalidates the speculative launch.
